\chapter{C++ and the STL for Contests}\label{ch:cpp}

C++ is the language of competitive programming for two reasons: it compiles to fast machine code, and its standard library
(the \defn{STL}) already contains most of the data structures and algorithms you need, implemented and debugged by experts.
Knowing the STL well is worth as much as knowing a dozen algorithms: every minute you do not spend writing and debugging a
balanced tree or a heap is a minute spent on the actual problem.

This chapter is a working reference, organised by what you need to do. It explains the parts of C++ you will use every day,
what each operation costs (measured, where it matters), and the traps that cost people contests.

\begin{chapterintro}{\bronze\ \silver\ \gold}{Chapter~\ref{ch:complexity}. The code assumes your template.}
  \item read and write data fast, and why the defaults are ten to fifty times slower;
  \item choose integer and floating-point types that cannot overflow;
  \item pick the right container for the job, and know the cost of every operation;
  \item use iterators, binary search, sorting and the rest of \texttt{<algorithm>} fluently;
  \item write recursive lambdas, use policy-based trees, random numbers, and avoid undefined behaviour.
\end{chapterintro}

% ======================================================================================================================
\section[Input and output]{Input and output \hfill \bronze}\label{sec:io}

By default, C++ streams stay synchronised with C's \cc{printf} and \cc{scanf} so that you may mix the two, and \cc{cin}
flushes \cc{cout} before every read so that prompts appear before input is requested. Both courtesies are expensive, and
contests need neither. Your template turns them off:

\begin{code}
ios::sync_with_stdio(false);  // stop synchronising with scanf / printf
cin.tie(nullptr);             // do not flush cout before every read
\end{code}

The difference is dramatic. Reading $10^6$ integers takes 0.72\,s with the default settings and 0.08\,s without
synchronisation; even \cc{scanf} needs 0.27\,s. For output, printing $10^6$ lines with \cc{endl}, which flushes the buffer
every time, takes almost 2 seconds, while \cc|'\n'| with synchronisation off takes 0.04\,s: forty-seven times faster. Every
one of those numbers was measured on the same machine as Table~\ref{tab:measured}.

\begin{pitfall}[After turning sync off]
Never mix \cc{cin}/\cc{cout} with \cc{scanf}/\cc{printf} once synchronisation is off: the two sides buffer separately and
their output comes out interleaved in the wrong order. Pick one family per program.
\end{pitfall}

\subsection{Reading patterns}

\begin{code}
int n;
cin >> n;
vi a(n);
cin >> a;                     // template: reads a whole vector
string s;
cin >> s;                     // one word (stops at whitespace)
string line;
getline(cin >> ws, line);     // a whole line (ws skips the old newline)
ll x;
while (cin >> x) { /* ... */ }  // until the end of input (count not given)
vector<string> grid(r);       // a grid of characters, one row per word
cin >> grid;
\end{code}

The one trap in reading is mixing \cc{>>} with \cc{getline}: after \cc{cin >> n}, the newline that ended the line is still
in the input, and a plain \cc{getline} returns an empty string. Skipping whitespace with \cc{cin >> ws} first, as above, fixes
it.

\subsection{Writing patterns}

\begin{code}
print(a, b, c);               // template: "a b c\n"
print(v);                     // template: a vector, space separated
yn(ok);                       // template: YES / NO
cout << fixed << setprecision(10) << ans << '\n';  // the template sets this
cout << flush;                // interactive problems: flush after each query
\end{code}

Floating-point answers are usually checked with an absolute or relative error of $10^{-6}$, so print more digits than
that: the template's 10 digits after the decimal point are safe. For interactive problems (Chapter~\ref{ch:special}), the
judge only sees your query after you flush, so there \cc{endl} or an explicit flush is required, not a mistake.

% ======================================================================================================================
\section[Integer types and overflow]{Integer types and overflow \hfill \bronze}\label{sec:overflow}

Overflow is the most common cause of wrong answers that pass the samples. A signed integer that overflows is
\emph{undefined behaviour} in C++: in practice it wraps around to a negative number, but the compiler is also allowed to
assume it never happens and optimise accordingly, which can make the bug appear or vanish between compilers. The defence is to
know your ranges (Figure~\ref{fig:ranges}) and to estimate the largest value every expression can reach.

\begin{figure}[htbp]
\centering
\begin{tikzpicture}[x=0.3cm]
  % a logarithmic ruler: position = log10(value)
  \draw[->, line width=0.7pt, fgink!70] (0, 0) -- (40.8, 0) node[right, lbl] {value};
  \foreach \e in {0, 5, 10, 15, 20, 25, 30, 35, 40} {
    \draw[line width=0.5pt, fgink!70] (\e, 0) -- (\e, 0.1);
    \node[font=\sffamily\tiny, text=fggray, anchor=south, inner sep=1pt] at (\e, 0.1) {$10^{\e}$};
  }
  % type ranges as bars above the ruler (they start above the tick labels)
  \foreach \v/\y/\c/\t in {9.33/0.75/fgblue/{\texttt{int}: $2.1\cdot10^9$},
      15.95/1.3/fgorange/{\texttt{double} (exact integers): $2^{53} \approx 9\cdot10^{15}$},
      18.96/1.85/fggreen/{\texttt{long long}: $9.2\cdot10^{18}$}, 19.27/2.4/fgpink/{\texttt{unsigned long long}: $1.8\cdot10^{19}$},
      38.23/2.95/fgsky/{\texttt{\_\_int128}: $1.7\cdot10^{38}$}} {
    \fill[fill=\c!35, draw=\c!80!black, line width=0.4pt] (0, \y-0.16) rectangle (\v, \y+0.16);
    \node[font=\sffamily\scriptsize, anchor=west] at (\v+0.3, \y) {\t};
  }
  % typical contest quantities: numbered markers on the ruler, explained underneath
  \foreach \v/\k in {5.3/1, 10.6/2, 14.3/3, 18.0/4} {
    \fill[fgverm] (\v, 0) circle (0.9mm);
    \node[circle, fill=fgverm, text=white, font=\sffamily\tiny\bfseries, inner sep=0.8pt, minimum size=3.6mm] at (\v, -0.95) {\k};
    \draw[line width=0.6pt, fgverm] (\v, -0.05) -- (\v, -0.77);
  }
  \node[font=\sffamily\scriptsize, text=fgink, anchor=north west, text width=12.4cm] at (0, -1.4) {
    \textcolor{fgverm}{\textbf{1}} $n = 2\cdot10^5$ elements \quad
    \textcolor{fgverm}{\textbf{2}} $n^2 = 4\cdot10^{10}$ pairs \quad
    \textcolor{fgverm}{\textbf{3}} a sum of $n$ values up to $10^9$: $2\cdot10^{14}$ \quad
    \textcolor{fgverm}{\textbf{4}} a product of two values up to $10^9$: $10^{18}$};
\end{tikzpicture}
\caption{Type ranges (bars) against quantities that appear in contests (red ticks), on a logarithmic ruler. A sum of $2\cdot
10^5$ values up to $10^9$ overflows \texttt{int} by a factor of $10^5$; a product of two such values fits in \texttt{long long}
with a factor of 9 to spare, but a product of three does not fit anywhere except \texttt{\_\_int128}.}
\label{fig:ranges}
\end{figure}

\begin{pitfall}[Overflow happens in the expression, not the variable]
\cc{long long x = a * b;} still overflows when \cc{a} and \cc{b} are 32-bit \cc{int}s, because the product is computed in
\cc{int} before it is stored. Write \cc{1LL * a * b} or make the operands 64-bit. Your template's \cc{#define int long long}
removes most of these bugs, but not the ones in \cc{1 << k} (the literal \cc{1} is still a 32-bit \cc{int}: write
\cc{1LL << k}), in \cc{__builtin_popcount} (use \cc{__builtin_popcountll}), or in constants like \cc{1e18}, which is a
\cc{double}.
\end{pitfall}

\paragraph{When 64 bits are not enough.} A product of two numbers below $9.2 \cdot 10^{18}$ can reach $8.5 \cdot 10^{37}$,
which fits in GCC's \cc{__int128} (but not on 32-bit judges, and \cc{cin}/\cc{cout} cannot read or print it directly). For
modular multiplication with a modulus up to $10^{18}$, the portable alternatives are a \cc{long double} trick or splitting
into 32-bit halves, both in your library (\texttt{cpmodmul} ideas inside \texttt{cppollard} and \texttt{cphash}). To
\emph{detect} overflow instead, GCC provides \cc{__builtin_mul_overflow(a, b, &res)}, which returns true when the product
does not fit.

\paragraph{Unsigned arithmetic wraps by design.} \cc{unsigned long long} arithmetic is defined modulo $2^{64}$, which makes
it perfect for hashing (Chapter~\ref{ch:strings}) and dangerous for everything else: \cc{v.size() - 1} is a huge number when
the vector is empty, because \cc{size()} is unsigned. The template's \cc{sz(v)} returns a signed value for exactly this
reason.

\paragraph{Floating point.} A \cc{double} has 53 bits of mantissa: about 15--16 significant decimal digits, and every
integer up to $2^{53} \approx 9 \cdot 10^{15}$ is exact. \cc{long double} on x86 GCC has 64 bits of mantissa (18--19 digits).
Floating-point answers are fine for geometry and probability; for anything that must be exact (counts, sums of large
integers, comparisons of fractions), use integers and cross-multiplication.

\begin{fact}[Powers of two to remember]
$2^{10} \approx 10^3$, $2^{20} \approx 10^6$, $2^{30} \approx 10^9$, $2^{31} - 1 = 2\,147\,483\,647$, $2^{53} \approx 9 \cdot
10^{15}$, $2^{60} \approx 10^{18}$, $2^{63} - 1 \approx 9.2\cdot10^{18}$. A number up to $10^{18}$ has 60 bits, so binary
search over $[0, 10^{18}]$ takes 60 steps.
\end{fact}

\begin{pitfall}[Division and remainder with negative numbers]
C++ rounds integer division towards zero: $-7 / 2 = -3$ and $-7 \bmod 2 = -1$. When you need the mathematical floor or a
non-negative remainder, use the template's \cc{fdiv(a, b)} and \cc{cdiv(a, b)}, and write \cc{((x % m) + m) % m}.
\end{pitfall}

% ======================================================================================================================
\section[Choosing a container]{Choosing a container \hfill \silver}\label{sec:containers}

The STL has a container for almost every access pattern, and choosing the right one is most of the work. Figure
\ref{fig:container-choice} is a decision procedure; Table~\ref{tab:containers} lists what each operation costs.

\begin{figure}[htbp]
\centering
\begin{tikzpicture}[q/.style={fq=6.2cm, minimum height=0.9cm}, c/.style={fend=3.1cm, minimum height=0.9cm}]
  \foreach \t/\ans [count=\k] in {
      {Only append, and access by position?}/{vector},
      {Add or remove at both ends?}/{deque},
      {Repeatedly take the smallest or largest?}/{priority\_queue},
      {Keep items sorted; find the next larger, or the neighbours of $x$?}/{set / multiset},
      {Need the $k$-th smallest, or the rank of $x$?}/{ordered\_set (PBDS)},
      {Look up values by key, order irrelevant?}/{unordered\_map},
      {Look up by key, and iterate in key order?}/{map},
      {Fixed-size boolean array with bulk operations?}/{bitset}} {
    \pgfmathsetmacro{\y}{-1.18*(\k-1)}
    \node[q] (q\k) at (0, \y) {\t};
    \node[c] (c\k) at (7.2, \y) {\texttt{\ans}};
    \draw[farr] (q\k) -- node[fyes, above] {yes} (c\k);
  }
  \foreach \k [evaluate=\k as \nk using int(\k+1)] in {1,...,7} \draw[farr] (q\k) -- node[fno, right] {no} (q\nk);
\end{tikzpicture}
\caption{Choosing a container: ask the questions from the top and take the first ``yes''. Most problems stop at the first
question; \texttt{vector} is the right answer far more often than any other container.}
\label{fig:container-choice}
\end{figure}

\begin{table}[htbp]
\centering
\caption{STL containers and the cost of their operations ($n$ = number of elements).}\label{tab:containers}
\small
\begin{tabularx}{\linewidth}{P{3.1cm}P{2.9cm}P{3.3cm}L}
\toprule
\thd{Container} & \thd{Access / search} & \thd{Insert / erase} & \thd{Typical use} \\
\midrule
\texttt{vector} & $\Oh{1}$ by index & $\Oh{1}$ amortized at the end & almost everything \\
\texttt{deque} & $\Oh{1}$ by index & $\Oh{1}$ at both ends & sliding windows, 0-1 BFS \\
\texttt{stack}, \texttt{queue} & top / front only & $\Oh{1}$ & DFS, BFS, bracket matching \\
\texttt{priority\_queue} & top in $\Oh{1}$ & $\Oh{\log n}$ & Dijkstra, ``take the best'' \\
\texttt{set}, \texttt{map} & $\Oh{\log n}$ & $\Oh{\log n}$ & sorted keys, neighbours \\
\texttt{multiset} & $\Oh{\log n}$ & $\Oh{\log n}$ & sorted with duplicates, medians \\
\texttt{unordered\_map} & $\Oh{1}$ expected & $\Oh{1}$ expected & counting by key \\
\texttt{bitset<N>} & $\Oh{1}$ & $\Oh{N/64}$ per bulk op & boolean DP, reachability \\
\texttt{string} & $\Oh{1}$ by index & $\Oh{1}$ at the end & text; \texttt{substr} copies \\
\bottomrule
\end{tabularx}
\end{table}

\subsection{vector, and resetting between test cases}

\begin{code}
vi v(n, 0);                            // n zeros
vvi g(n, vi(m, -1));                   // an n x m grid of -1
v.push_back(x); v.pop_back(); v.back(); sz(v); v.empty();
v.assign(n, 0);                        // reset to n zeros, reusing the memory
v.clear();                             // size 0, capacity kept
sort(all(v)); reverse(all(v));
v.erase(unique(all(v)), v.end());      // remove duplicates (sort first)
v.reserve(n);                          // no reallocations if the size is known
\end{code}

\cc{vector} stores its elements contiguously, which makes it the fastest container to scan by a wide margin
(Section~\ref{sec:measured}). In multi-test problems, declare vectors inside \cc{solve()} so that each test starts fresh,
or reset them with \cc{assign}: a global array cleared with \cc{memset} in every test costs its full size each time, even
when the test is tiny.

\subsection{Iterators and half-open ranges}

Every STL algorithm works on a \defn{half-open range} $[\mathit{first}, \mathit{last})$: \cc{begin()} points at the first
element and \cc{end()} points one past the last (Figure~\ref{fig:iterators}). Half-open ranges compose without off-by-one
errors: $[a, b)$ followed by $[b, c)$ is $[a, c)$, the length is $\mathit{last} - \mathit{first}$, and an empty range has
$\mathit{first} = \mathit{last}$. The whole of this book uses the same convention for array ranges.

\begin{figure}[htbp]
\centering
\begin{tikzpicture}
  \arr{A}{0}{0}{1,3,3,3,5,8,9}
  \node[cell, dim, dashed] (A-8) at ($(A-7.east)$) {};
  \node[lbl, anchor=east] at (-0.2, 0) {sorted \texttt{v}};
  % above: the two bounds for x = 3; below: begin(), end() and the run of 3s
  \ptra{A-2}{lower\_bound}{fgverm}
  \ptra{A-5}{upper\_bound}{fggreen!80!black}
  \ptrb{A-1}{begin()}{fgblue}
  \ptrb{A-8}{end()}{fgblue}
  \draw[mbrace] ($(A-2.south west)+(0,-0.08)$) -- ($(A-4.south east)+(0,-0.08)$);
  \node[font=\sffamily\scriptsize, text=fgpink!80!black] at ($(A-3.south)+(0,-0.72)$) {the three 3s};
  \node[notel, anchor=west, text width=5.2cm] at (6.3, -0.35) {\texttt{lower\_bound(x)}: first element $\ge x$.\\
    \texttt{upper\_bound(x)}: first element $> x$.\\ Between them lie exactly the copies of $x$: here 3 of them, so
    \texttt{count = upper $-$ lower}. \texttt{end()} is one past the last element and must never be dereferenced.};
\end{tikzpicture}
\caption{Iterators on a sorted vector. Both bounds take $\Oh{\log n}$ and return \texttt{end()} when no element qualifies.}
\label{fig:iterators}
\end{figure}

\begin{code}
auto it = lower_bound(all(v), x);      // first element >= x (v sorted)
int idx = it - v.begin();              // its index; idx == sz(v) means "none"
auto [lo, hi] = equal_range(all(v), x); // both bounds at once
bool present = binary_search(all(v), x);
int cnt = upper_bound(all(v), x) - lower_bound(all(v), x);
\end{code}

\subsection{set, multiset and map}

\cc{set} and \cc{map} are balanced binary search trees (red-black trees in GCC). Every operation walks one root-to-leaf path,
so it costs $\Oh{\log n}$, and the elements are always available in sorted order. Their power lies in the neighbour queries:
\cc{lower_bound} walks down the tree remembering the smallest key that was not too small (Figure~\ref{fig:bst-lower}).

\begin{figure}[htbp]
\centering
\begin{forest}
  cptree, for tree={s sep=3.5mm, l sep=5mm, minimum size=6.8mm}
  [23, hot
    [11, hot
      [5 [2] [8]]
      [17, good
        [14, hot] [20]]]
    [35
      [29 [26] [32]]
      [41 [38] [44]]]]
\end{forest}
\par\medskip
{\small
\begin{tabular}{@{}llllll@{}}
\multicolumn{6}{@{}l}{\texttt{s.lower\_bound(15)} visits:} \\[2pt]
\toprule
node & 23 & 11 & 17 & 14 & \\
compare & $15 < 23$: left & $15 > 11$: right & $15 < 17$: left & $15 > 14$: right & \\
best so far & 23 & 23 & \textbf{17} & 17 & $\Rightarrow$ answer 17 \\
\bottomrule
\end{tabular}}
\caption{\texttt{lower\_bound} in a balanced search tree walks one path from the root (orange), remembering the last node
where it went left. With $n$ keys the path has about $\log_2 n$ nodes.}
\label{fig:bst-lower}
\end{figure}

\begin{code}
set<int> s;
s.insert(x); s.erase(x); s.count(x);   // count is 0 or 1
auto it = s.lower_bound(x);            // first element >= x, or s.end()
if (it != s.begin()) { int below = *prev(it); }  // largest element < x
*s.begin(); *s.rbegin();               // minimum, maximum
multiset<int> ms;
ms.erase(ms.find(x));                  // erase ONE copy (erase(x) erases all!)
map<string, int> cnt;
cnt[w]++;                              // a missing key starts at 0
for (auto& [key, value] : cnt) { }     // in increasing key order
\end{code}

\begin{pitfall}[Use the member function]
\cc{lower_bound(all(s), x)} on a \cc{set} compiles, but it is the generic algorithm, which moves an iterator step by step
and costs $\Oh{n}$. Always call \cc{s.lower_bound(x)}. The same goes for \cc{s.count(x)} versus \cc{count(all(s), x)}.
\end{pitfall}

\subsection{priority\_queue and the binary heap}

A \cc{priority_queue} is a \defn{binary heap} stored in a vector: the element at index $i$ has its children at $2i + 1$
and $2i + 2$, and every parent is at least as large as its children (Figure~\ref{fig:heap}). The largest element sits at the
root, \cc{push} lets a new element rise to its place and \cc{pop} lets the last element sink from the root, each along one
path of length $\log_2 n$.

\begin{figure}[htbp]
\centering
\begin{tikzpicture}
  \begin{scope}
    \node[vtx, hot] (h0) at (0, 0) {9};
    \node[vtx] (h1) at (-1.8, -1.05) {7}; \node[vtx] (h2) at (1.8, -1.05) {8};
    \node[vtx] (h3) at (-2.7, -2.1) {3}; \node[vtx] (h4) at (-0.9, -2.1) {6};
    \node[vtx] (h5) at (0.9, -2.1) {2}; \node[vtx] (h6) at (2.7, -2.1) {5};
    \node[vtx] (h7) at (-3.15, -3.15) {1};
    \draw[ed] (h0) -- (h1) (h0) -- (h2) (h1) -- (h3) (h1) -- (h4) (h2) -- (h5) (h2) -- (h6) (h3) -- (h7);
    \foreach \i in {0,...,7} \node[font=\sffamily\tiny, text=fggray, anchor=east] at ($(h\i.west)+(-0.02,0)$) {\i};
  \end{scope}
  \begin{scope}[xshift=4.6cm, yshift=-1.3cm]
    \arrs{H}{0}{0}{9/hot, 7/n, 8/n, 3/n, 6/n, 2/n, 5/n, 1/n}
    \idx{H}{8}
    \draw[->, line width=0.8pt, fgblue] (H-2.north) to[out=60, in=120] (H-4.north);
    \draw[->, line width=0.8pt, fgblue] (H-2.north) to[out=55, in=125] (H-5.north);
    \node[font=\sffamily\scriptsize, text=fgblue] at (2.1, 1.3) {children of $i$: $2i + 1$, $2i + 2$};
  \end{scope}
\end{tikzpicture}
\caption{A max-heap drawn as a tree (left) and stored as an array (right). The tree is complete, so the array has no gaps,
and no pointers are needed: the parent of index $i$ is $\lfloor (i - 1)/2 \rfloor$.}
\label{fig:heap}
\end{figure}

\begin{code}
priority_queue<int> pq;                // max-heap: top() is the largest
minpq<pll> q;                          // min-heap of (distance, vertex)
q.push({d, v});
auto [dist, u] = q.top(); q.pop();     // pairs compare by first, then by second
\end{code}

A heap cannot find or delete an arbitrary element. When you need that, either use a \cc{set} (which can) or use
\emph{lazy deletion}: leave outdated entries in the heap and skip them when they surface, as Dijkstra's algorithm does
(Chapter~\ref{ch:shortest}).

\subsection{unordered\_map and hash attacks}

\cc{unordered_map} is a hash table: an array of buckets, each holding the keys whose hash falls into it. With a good hash,
buckets are short and every operation is expected $\Oh{1}$. GCC's hash of an integer is the integer itself, and the bucket is
the hash modulo a prime from a fixed list, so an adversary who knows that prime can send keys that all land in one bucket
(Figure~\ref{fig:hashattack}). Every operation then scans the whole chain: $\Oh{n}$ each, $\Oh{n^2}$ in total. On Codeforces,
where participants hack each other's solutions, this happens routinely.

\begin{figure}[htbp]
\centering
\begin{tikzpicture}[b/.style={draw=fgink!60, fill=fgpaper, minimum width=6mm, minimum height=5mm, inner sep=0pt,
    font=\sffamily\scriptsize}, k/.style={draw=fgverm, fill=fgverm!15, rounded corners=1pt, minimum height=4.5mm, inner sep=2pt,
    font=\ttfamily\scriptsize}, kg/.style={k, draw=fggreen, fill=fggreen!15}]
  \node[font=\sffamily\small\bfseries, anchor=west] at (-0.3, 0.75) {keys that are multiples of the bucket count};
  \foreach \i in {0,...,6} \node[b] (A\i) at (\i*0.65, 0) {\i};
  \node[k] (a1) at (0, -0.8) {0}; \node[k] (a2) at (0, -1.4) {97}; \node[k] (a3) at (0, -2.0) {194};
  \node[k] (a4) at (0, -2.6) {291}; \node[font=\small, text=fgverm] at (0, -3.1) {$\vdots$};
  \draw[line width=0.6pt, fgverm] (A0) -- (a1) -- (a2) -- (a3) -- (a4);
  \node[lbl, text width=3.8cm, align=left, anchor=west] at (1.1, -1.9) {every key lands in bucket 0: each lookup scans
    the whole chain, $\Oh{n}$};
  \begin{scope}[xshift=7.2cm]
    \node[font=\sffamily\small\bfseries, anchor=west] at (-0.3, 0.75) {the same keys, randomised hash};
    \foreach \i in {0,...,6} \node[b] (B\i) at (\i*0.65, 0) {\i};
    \foreach \i/\key in {0/291, 2/0, 3/194, 5/97} {
      \node[kg] (k\i) at (\i*0.65, -0.8) {\key};
      \draw[line width=0.6pt, fggreen] (B\i) -- (k\i);
    }
    \node[lbl, text width=4.3cm, align=left, anchor=west] at (-0.3, -1.9) {a random seed inside the hash spreads them out:
      expected $\Oh{1}$, and no one can predict the seed};
  \end{scope}
\end{tikzpicture}
\caption{An anti-hash test against \texttt{unordered\_map} (left) and the fix (right): hash with a random seed, as the
library's \texttt{custom\_hash} does, so that no fixed input can force collisions.}
\label{fig:hashattack}
\end{figure}

\begin{code}
unordered_map<ll, int, custom_hash> cnt;  // template cphash: seeded splitmix64
cnt.reserve(1 << 20);                     // pre-size the buckets: 1.6x faster
\end{code}

\subsection{bitset}

\cc{bitset<N>} is a fixed-size array of $N$ bits packed into 64-bit words. Single bits cost $\Oh{1}$; whole-set operations
(\cc{&}, \cc{|}, \cc{^}, \cc{~}, shifts, \cc{count()}, \cc{any()}) cost $\Oh{N/64}$ because they process 64 bits per machine
instruction. GCC adds \cc{_Find_first()} and \cc{_Find_next(i)} to jump between set bits. The size must be a compile-time
constant; when it depends on the input, round up to the largest possible value.

\begin{code}
bitset<100001> can;                     // subset sums reachable so far
can[0] = 1;
for (int w : weights) can |= can << w;  // N/64 word operations per item
int count_reachable = can.count();
for (int i = can._Find_first(); i < 100001; i = can._Find_next(i)) {
    // i is a reachable sum
}
\end{code}

% ======================================================================================================================
\section[Sorting and comparators]{Sorting and comparators \hfill \silver}\label{sec:sorting-stl}

\cc{std::sort} is an introsort: quicksort that falls back to heapsort on bad inputs, so it is $\Oh{n\log n}$ in the worst
case, and it sorts $10^6$ integers in about 0.04\,s. It sorts by \cc{operator<} unless you pass a comparator.

\begin{code}
sort(all(v));                                    // increasing
sort(all(v), greater<>());                       // decreasing
sort(all(p), [](const pll& a, const pll& b) {    // by second, then first desc.
    if (a.second != b.second) return a.second < b.second;
    return a.first > b.first;
});
stable_sort(all(v), cmp);                        // equal elements keep order
vi ord(n);                                       // sort indices by a key
iota(all(ord), 0);
sort(all(ord), [&](int i, int j) { return key[i] < key[j]; });
\end{code}

\begin{pitfall}[Strict weak ordering]
A comparator must behave like \cc{<}: \cc{cmp(x, x)} must be false, and if \cc{cmp(a, b)} and \cc{cmp(b, c)} then
\cc{cmp(a, c)}. Writing \cc{<=} instead of \cc{<} breaks the first rule, and \cc{sort} may then read past the end of the
array and crash. Your local builds use the checked STL, which detects an invalid comparator and stops with a message.
\end{pitfall}

\paragraph{Sorting pairs and tuples.} Pairs, tuples and arrays compare lexicographically: first by the first component, ties
by the second, and so on. That makes ``sort by $a$, then by $b$ decreasing'' a one-liner: sort the pairs
\cc|{a, -b}|. Sorting
events as \cc{(time, type)} pairs orders simultaneous events by type for free, which Section~\ref{sec:events} used.

\paragraph{Coordinate compression.} Values up to $10^9$ that only matter through their order can be replaced by their rank
in $0, \dots, k - 1$, so they can index an array or a Fenwick tree (template \texttt{cpcompress}):

\begin{code}
vi vals = a;
sort(all(vals)), vals.erase(unique(all(vals)), vals.end());
for (auto& x : a) x = lower_bound(all(vals), x) - vals.begin();  // 0 <= x < sz
\end{code}

% ======================================================================================================================
\section[The algorithm toolbox]{The algorithm toolbox \hfill \silver}\label{sec:toolbox}

\begin{center}
\small
\begin{tabularx}{\linewidth}{>{\ttfamily\raggedright\arraybackslash}p{6.3cm}L}
\toprule
\thd{Call} & \thd{What it does} \\
\midrule
accumulate(all(v), 0LL) & sum; the \texttt{0LL} makes the sum 64-bit \\
*max\_element(all(v)), min\_element & largest / smallest element (an iterator) \\
count(all(v), x), find(all(v), x) & occurrences / first occurrence: $\Oh{n}$ \\
iota(all(v), 0) & fills $0, 1, 2, \dots$ \\
partial\_sum, adjacent\_difference & prefix sums and differences, in place \\
next\_permutation(all(v)) & next arrangement in lexicographic order; \texttt{false} after the last \\
nth\_element(v.begin() + k, \ldots) & the $k$-th smallest in place, expected $\Oh{n}$ \\
rotate, reverse, unique, fill & rearrangements, all $\Oh{n}$ \\
merge, inplace\_merge & merge two sorted ranges in $\Oh{n}$ \\
gcd(a, b), lcm(a, b) & greatest common divisor, least common multiple \\
\_\_builtin\_popcountll(x) & number of set bits \\
\_\_builtin\_ctzll(x), \_\_builtin\_clzll(x) & trailing / leading zero bits ($x > 0$) \\
\_\_lg(x) & $\lfloor \log_2 x \rfloor$ for $x > 0$ \\
shuffle(all(v), rng) & random order \\
\bottomrule
\end{tabularx}
\end{center}

\begin{trick}[nth\_element for medians and top-$k$]
\cc{nth_element(v.begin(), v.begin() + k, v.end())} puts the element that would be at index $k$ after sorting into position
$k$, with smaller elements before it and larger after, in expected linear time. It finds a median in one call and is several
times faster than sorting when you only need the $k$ smallest elements.
\end{trick}

% ======================================================================================================================
\section[Modern C++ that saves time]{Modern C++ that saves time \hfill \silver}\label{sec:modern}

C++17 has several features that make contest code shorter and less error-prone. All of them are used throughout this book.

\begin{code}
for (auto& x : v) x *= 2;              // by reference: modifies the elements
for (const auto& s : words) { }        // const& avoids copying strings
auto [a, b] = pr;                      // structured bindings unpack pairs etc.
for (auto& [key, val] : mp) { }
array<int, 3> t = {1, 2, 3};           // fixed size, usable as a map key
v.emplace_back(a, b);                  // construct a pair or struct in place
if (auto it = mp.find(k); it != mp.end()) { }  // "it" stays local to the if
\end{code}

\subsection{Lambdas and recursion}

A lambda is a function defined inside another function. \cc{[&]} lets it read and modify every local variable by reference,
which is why lambdas are the natural way to write helpers inside \cc{solve()}. A lambda cannot call itself by name, so a
recursive lambda receives itself as its first parameter:

\begin{code}
auto dfs = [&](auto&& self, int u, int p) -> void {
    for (int v : adj[u])
        if (v != p) self(self, v, u);
};
dfs(dfs, 0, -1);
\end{code}

This style is used by the library's inside-\cc{solve()} templates: no global arrays that must be reset between test cases.
It is also faster than \cc{std::function}, which adds an indirect call and may allocate.

% ======================================================================================================================
\section[Policy-based data structures]{Policy-based data structures \hfill \gold}\label{sec:pbds}

GCC ships an extension library with a balanced tree that also knows the \emph{size of every subtree}, so it can answer
``what is the $k$-th smallest element?'' and ``how many elements are smaller than $x$?'' in $\Oh{\log n}$. \cc{std::set}
cannot do either.

\begin{code}
ordered_set<int> s;                    // template cposet
s.insert(5); s.insert(1); s.insert(9);
*s.find_by_order(1);                   // 5: the element at index 1 (sorted)
s.order_of_key(6);                     // 2: the number of elements < 6
\end{code}

The tree is a set, so duplicates are dropped; for a multiset, store pairs \cc|{value, unique_id}|. It is about twice as slow
as \cc{std::set}, which is still fast enough for $2 \cdot 10^5$ operations, and it replaces a Fenwick tree over compressed
values in many problems (Josephus Problem II, sliding medians, counting inversions online).

% ======================================================================================================================
\section[Randomness]{Randomness \hfill \gold}\label{sec:random-stl}

\cc{rand()} is slow, returns only 15 random bits on some compilers, and is predictable. Use a 64-bit Mersenne Twister seeded
from the clock (template \texttt{cprng}):

\begin{code}
mt19937_64 rng(chrono::steady_clock::now().time_since_epoch().count());
ll r = uniform_int_distribution<ll>(lo, hi)(rng);   // uniform in [lo, hi]
shuffle(all(v), rng);
\end{code}

Seeding from the clock matters on Codeforces, where other contestants can read your code: with a fixed seed, they can
construct an input that makes your randomised algorithm take its worst case. Chapter~\ref{ch:special} shows what randomness
buys: hashing, randomised pivots, random restarts and probabilistic checks.

% ======================================================================================================================
\section[Undefined behaviour and common bugs]{Undefined behaviour and common bugs \hfill \bronze}\label{sec:ub}

Some mistakes in C++ are not errors the compiler reports but \defn{undefined behaviour}: the program may crash, print
garbage, or work perfectly on your machine and fail on the judge. The checked STL and the flags in your setup catch many of
them locally (GUIDE.pdf, Chapter~5); the rest you must avoid by habit.

\begin{itemize}
  \item \textbf{Out-of-bounds access}: \cc{a[n]} in an array of size $n$, \cc{v[-1]}, \cc{*s.begin()} on an empty set. The
        checked STL catches it for \cc{vector} and \cc{string} and names the line.
  \item \textbf{Uninitialised variables}: a local \cc{int sum;} holds garbage, not 0. Global variables start at 0.
  \item \textbf{Signed overflow}: undefined, not merely wrapping (Section~\ref{sec:overflow}).
  \item \textbf{Invalid iterators}: an iterator or reference into a \cc{vector} dies when the vector reallocates
        (\cc{push_back}); erasing from a container while iterating over it with range-for skips or repeats elements.
  \item \textbf{Missing return}: a non-\cc{void} function that ends without \cc{return}. With optimisation, GCC may turn
        this into a crash or an infinite loop; \texttt{-Wall} warns.
  \item \textbf{Bad comparator}: see Section~\ref{sec:sorting-stl}.
\end{itemize}

And the bugs that are not undefined but equally common: \cc{=} instead of \cc{==} in a condition (\texttt{-Wall} warns),
integer division where a real number was meant (\cc{1 / 2 == 0}), comparing doubles with \cc{==} instead of
\cc{fabs(a - b) < 1e-9}, forgetting to reset state between test cases, and reading the input in a different order from the
statement.

\begin{tip}[Debug output that cannot break the judge]
Your template's \cc{dbg(x, v, mp)} prints variables with their names to standard error locally, and compiles to nothing on
the judge. Leave the calls in: they cost nothing when submitted, and they are there when the next bug appears.
\end{tip}

\begin{summary}[Chapter summary]
\begin{itemize}
  \item Fast I/O (sync off, \cc|'\n'|) is worth a factor of 10 to 50; the template sets it up.
  \item Estimate every expression's largest value; use 64-bit integers, \cc{1LL <<}, and exact integer arithmetic.
  \item \cc{vector} first; \cc{set}/\cc{map} for sorted neighbours; \cc{priority_queue} for repeated minimums;
        \cc{unordered_map} with a randomised hash for counting; \cc{bitset} for bulk boolean work.
  \item Ranges are half-open; call member functions like \cc{s.lower_bound}; comparators must be strict.
  \item Lambdas with \cc{[&]}, structured bindings and \cc{const&} make contest code short and safe.
\end{itemize}
\end{summary}

\begin{practice}
\ppgroup{Containers in action}
\begin{pplist}
  \item Distinct Numbers
  \item Apartments
  \item Concert Tickets
  \item Sum of Two Values
  \item Playlist
  \item Towers
  \item Traffic Lights
  \item Room Allocation
  \item Collecting Numbers II
\end{pplist}
\ppgroup{Ordered statistics}
\begin{pplist}
  \item Josephus Problem II
  \item Sliding Median
  \item Sliding Cost
\end{pplist}
\end{practice}
